%-----------------------------------------------------------
%   Capítulo 1 - Introdução
%-----------------------------------------------------------
\chapter{Introdução}
\label{cap:introducao}

\section{Enquadramento e Motivação}

O projeto de um sistema de leitura e escrita multiprotocolo insere-se
no contexto dos sistemas embebidos e autónomos em que coexistem diferentes
tecnologias de comunicação. O trabalho abrange I2C, RS232, RS485, CAN e
Ethernet e é apresentado neste relatório numa fase preliminar de definição,
centrada no enquadramento do problema e na organização do estudo e do
desenvolvimento.

Os sistemas autónomos e os sistemas embebidos integram funções de
perceção, processamento, controlo e atuação que dependem da troca de
informação entre os seus elementos. A aquisição de dados de sensores,
o envio de comandos a atuadores e a comunicação com controladores ou
equipamentos de supervisão podem envolver tecnologias de comunicação
distintas numa mesma instalação. Esta coexistência resulta das
características dos dispositivos e das funções que desempenham, bem
como da necessidade de integrar equipamentos desenvolvidos em momentos
ou contextos diferentes.

Neste contexto, sensores, atuadores, controladores e outros equipamentos
podem disponibilizar interfaces como I2C, RS232, RS485, CAN e Ethernet.
Estas tecnologias apresentam diferenças nas características elétricas,
na organização das ligações e nos mecanismos de comunicação que
suportam. Por conseguinte, a ligação de um novo equipamento exige a
análise da compatibilidade da interface e da forma como os dados são
transmitidos e recebidos, além da compreensão do seu significado para
a aplicação.

O problema central reside na heterogeneidade destas interfaces e na
dificuldade de integrar operações de leitura e escrita sobre equipamentos
com formas de acesso distintas. O suporte de cada tecnologia pode exigir
configurações, procedimentos de comunicação e mecanismos de tratamento
de erros específicos. Quando estas particularidades são tratadas de
forma dispersa na aplicação, a inclusão ou substituição de equipamentos
pode implicar alterações em várias partes do sistema, dificultando a
manutenção e a reutilização das funções de comunicação. Além disso, a
disponibilidade de uma ligação compatível não assegura, por si só, a
interpretação correta dos dados trocados.

A relevância do trabalho reside, assim, na possibilidade de organizar de
forma coerente os mecanismos de acesso aos equipamentos, facilitando a
integração e a evolução de sistemas com diferentes tecnologias de
comunicação.

\section{Cenários de aplicação}

\textit{TODO: definir os cenários de aplicação após validação do âmbito do projeto.}

\section{Objectivos}

\subsection{Objetivo Geral}

Estudar e desenvolver um sistema capaz de realizar operações de leitura
e escrita através das tecnologias de comunicação I2C, RS232, RS485,
CAN e Ethernet.

\subsection{Objetivos Específicos}

Para concretizar o objetivo geral, estabelecem-se os seguintes objetivos
específicos:

\begin{enumerate}
    \item Estudar as características de I2C, RS232, RS485, CAN e Ethernet
    relevantes para a realização de operações de comunicação.

    \item Analisar as diferenças e limitações das cinco tecnologias,
    identificando as suas implicações para a integração num mesmo sistema.

    \item Definir uma arquitetura de integração das cinco interfaces,
    identificando os blocos funcionais e as respetivas ligações.

    \item Implementar as funções de comunicação necessárias à transmissão
    e receção de dados através de cada uma das cinco interfaces.

    \item Validar experimentalmente as operações de transmissão e receção
    em cada interface, de acordo com os requisitos que venham a ser
    consolidados.
\end{enumerate}

\section{Estrutura do Relatório}

O Capítulo~\ref{cap:introducao} apresenta o enquadramento, a motivação
e os objetivos do trabalho. O Capítulo~\ref{cap:estado-arte} organiza
o Estado da Arte sobre as tecnologias abrangidas.
O Capítulo~\ref{cap:requisitos} distingue as capacidades previstas,
as propostas preliminares e as questões a esclarecer com os docentes.
Por fim, o Capítulo~\ref{cap:arquitectura} apresenta os blocos funcionais
abstratos da Arquitetura de Alto Nível, cuja concretização depende
do estudo e da consolidação dos requisitos.
